% ===========================================================================
\section{Notación}
\label{sec:notacion}
% ===========================================================================

Se usa la misma notación BNF-extendida del documento \emph{Reglas de Sintaxis de MiniJava} de la cátedra, con la misma convención tipográfica, para que las reglas de este
informe se puedan comparar renglón a renglón con las de aquél.

\begin{center}
\begin{tabular}{@{}ll@{}}
\toprule
\textbf{Elemento} & \textbf{Significado} \\
\midrule
\tl{terminal}            & símbolo terminal \\
\nt{Clase}               & símbolo no terminal (primera letra mayúscula) \\
\eps                     & la cadena vacía \\
\nt{X} \deriva{} $\alpha$              & una producción, con $\alpha$ una secuencia de terminales y no terminales \\
\nt{X} \deriva{} $\alpha$ \alt{} $\beta$ & abreviación de \nt{X} \deriva{} $\alpha$ y \nt{X} \deriva{} $\beta$ \\
$\prim{\beta}$           & conjunto de terminales que pueden aparecer primeros en $\beta$ \\
\bottomrule
\end{tabular}
\end{center}

El no terminal inicial es \nt{Inicial}.

\subsection{El criterio LL(1)}

Una gramática es LL(1) si, para todo no terminal $A$ con producciones
$A \to \beta_1 \mid \beta_2 \mid \dots \mid \beta_n$, se cumple para todo $i \neq j$:

\[ \prim{\beta_i} \cap \prim{\beta_j} = \varnothing \]

Sólo con esa condición un único token de anticipación alcanza para elegir la producción correcta
sin necesidad de \emph{backtracking}. Cuando la condición no se cumple, la gramática debe
adaptarse \textbf{sin cambiar el lenguaje que genera}, que es la restricción que gobierna todas
las transformaciones de este informe.

\subsection{Dos defectos distintos que conviene no confundir}
\label{sec:dos-defectos}

Buena parte del trabajo consistió en separar dos cosas que se parecen y no son lo mismo.

\begin{ambiguedad}{el lenguaje está mal descripto}
Existe una cadena con \textbf{dos árboles de derivación distintos}. No queda determinado qué
significa el programa. No siempre se puede arreglar reescribiendo: a veces hay que \emph{elegir}
cuál de los dos significados vale.
\end{ambiguedad}

\begin{conflicto}{el analizador no puede predecir}
Dos producciones del mismo no terminal comparten Primeros. La cadena sigue teniendo \textbf{un
solo} árbol; lo que falla es que el analizador no puede saber cuál producción aplicar mirando un
token. Se corrige factorizando o eliminando recursión, y el lenguaje generado \textbf{no cambia}.
\end{conflicto}

Todo par de producciones que comparte prefijo izquierdo produce un conflicto LL(1); no
todo par que comparte prefijo izquierdo es ambiguo. Si los dos lados derivan lenguajes disjuntos,
no hay ninguna cadena con dos árboles.

Un criterio suficiente y barato para probar esa disyunción es comparar los \textbf{rangos de
longitud derivable}: si no se solapan, los lenguajes son disjuntos y no puede haber ambigüedad.
Se aplica en la sección~\ref{sec:diagnostico} y decide varios de los casos.
